\Askhsh{37017}{2}
\begin{ylh}{1}
    \item 4.2. Τέμνουσα δύο ευθειών - Ευκλείδειο αίτημα
    \item 4.6. Άθροισμα γωνιών τριγώνου
    \item 5.2. Παραλληλόγραμμα
    \item 5.9. Μια ιδιότητα του ορθογώνιου τριγώνου.
\end{ylh}
% ===================================================================
%  Αρχείο: 37017.tex (Fragment)
%  Αυτόματη μετατροπή μέσω doc2latex.py
% ===================================================================

ΘΕΜΑ 2

Σε παραλληλόγραμμο $ΑΒ\Gamma\Delta$ είναι $\widehat{Β} = 120^{\circ}$ και $\Delta Ε \perp B\Gamma$. Έστω $EZ$ η διάμεσος του τριγώνου $\Delta E\Gamma$.

\begin{alist}
\item Να υπολογίσετε τις γωνίες $\widehat{Α}$ και $\widehat{\Gamma}$ του παραλληλογράμμου. \monades{8}
\item Αν $K$ είναι το μέσο της πλευράς $AB$, να αποδείξετε ότι $EZ = AK$. \monades{9}
\item Να υπολογίσετε τη γωνία $E\widehat{Z}\Gamma$. \monades{8}

\begin{center}
\begin{tikzpicture}[scale=1]
\tkzDefPoint(0,0){D}
\tkzDefPoint(4,0){C}
\tkzDefPointBy[rotation=center D angle 120](C) \tkzGetPoint{A_ray}
\tkzDefPoint(2,0){temp}
\tkzDefPointBy[rotation=center D angle 120](temp) \tkzGetPoint{A}
\tkzDefPointBy[translation=from D to C](A) \tkzGetPoint{B}

\tkzDefLine[orthogonal=through D](B,C) \tkzGetPoint{DE_dir}
\tkzInterLL(D,DE_dir)(B,C) \tkzGetPoint{E}
\tkzDefMidPoint(D,C) \tkzGetPoint{Z}
\tkzDefMidPoint(A,B) \tkzGetPoint{K}

\draw[l3] (A)--(B)--(C)--(D)--cycle;
\draw[l3] (D)--(E);
\draw[l3] (E)--(Z);

\tkzDrawPoints(A,B,C,D,E,Z,K)
\tkzLabelPoint[above](A){$A$}
\tkzLabelPoint[above](B){$B$}
\tkzLabelPoint[below](C){$\Gamma$}
\tkzLabelPoint[below](D){$\Delta$}
\tkzLabelPoint[right](E){$E$}
\tkzLabelPoint[below](Z){$Z$}
\tkzLabelPoint[above](K){$K$}

\tkzMarkRightAngle(D,E,C)
\tkzMarkAngle[size=0.5](A,B,C)
\tkzLabelAngle[pos=0.8](A,B,C){$120^{\circ}$}
\end{tikzpicture}
\end{center}
\end{alist}
